\pdfvariable trailerid {[<C3362026090300000000000000000000><C3362026090300000000000000000000>]}
\documentclass[10pt]{article}
\usepackage[margin=0.72in]{geometry}
\usepackage{amsmath,amssymb,amsthm,mathtools,microtype,hyperref}
\hypersetup{colorlinks=true,linkcolor=blue,urlcolor=blue}
\providecommand{\CRevisionRound}{2}
\newtheorem{theorem}{Theorem}
\newtheorem{lemma}[theorem]{Lemma}
\newtheorem{remark}[theorem]{Remark}
\newcommand{\one}{\mathrm{1}}
\newcommand{\EE}{\mathbb E}
\newcommand{\cH}{\mathcal H}
\setlength{\emergencystretch}{3em}
\title{Single-Peak Crow--Kimura Dynamics:\\
Complete Finite-Genome Spectrum and Exact Projective Gap}
\author{Route-A source-local certificate HCS-C336}
\date{3 September 2026}

\begin{document}
\maketitle
\begin{abstract}
For the continuous-time mutation--selection equation on binary strings of
length $L$, with symmetric point mutation and one master-sequence fitness
spike, we identify the nonlinear simplex flow exactly with a projectivized
positive linear semigroup.  Walsh decomposition shows that each Hamming
eigenvalue retains an explicit codimension-one multiplicity, while all
remaining eigenvalues are the roots of one degree-$(L+1)$ secular polynomial.
\ifnum\CRevisionRound>0
The secular roots are simple and strictly interlace the mutation poles.  The
two largest roots therefore give the exact projective spectral gap, with a
sharp generic convergence law.
\fi
\ifnum\CRevisionRound>1
Zero selection, zero mutation and the one-locus model are closed separately;
finite-length analyticity is not promoted to an infinite-genome error
threshold.  Exact computation is used only as an implementation receipt.
\fi
\end{abstract}

\section{Frozen mutation--selection owner}

Let $\cH_L=\mathbb R^{\{0,1\}^L}$, let $F_i$ flip bit $i$, and let $e_0$
denote the all-zero genotype.  For $L\geq1$ and $U,s>0$ define
\begin{equation}\label{eq:operator}
 M_L=\frac UL\sum_{i=1}^L(F_i-I),\qquad
 A_L=M_L+s e_0e_0^{\mathsf T}.
\end{equation}
Both matrices are real symmetric and $\one^{\mathsf T}M_L=0$.  On the
probability simplex the Crow--Kimura equation in this normalization is
\begin{equation}\label{eq:nonlinear}
 \dot p=A_Lp-(\one^{\mathsf T}A_Lp)p=A_Lp-sp_0p.
\end{equation}

\begin{theorem}[Exact projectivization and global limit]\label{thm:flow}
For every probability vector $p^{\rm in}$,
\begin{equation}\label{eq:quotient}
 p(t)=\frac{e^{tA_L}p^{\rm in}}
 {\one^{\mathsf T}e^{tA_L}p^{\rm in}}
\end{equation}
is the unique solution of~\eqref{eq:nonlinear}.  The denominator is positive,
the simplex is preserved, and $p(t)$ converges to the uniquely normalized
positive Perron eigenvector of $A_L$.
\end{theorem}

\begin{proof}
Write $q(t)=e^{tA_L}p^{\rm in}$ and $Z(t)=\one^{\mathsf T}q(t)$.  Positivity of
the irreducible hypercube semigroup gives $q(t)>0$ for $t>0$, hence $Z(t)>0$.
Moreover $Z'=\one^{\mathsf T}A_Lq=sq_0$.  Differentiating $q/Z$ gives exactly
\eqref{eq:nonlinear}.  The same calculation shows that its entries sum to
one.  Since $A_L$ is irreducible Metzler and symmetric, its largest eigenvalue
is simple and has a strictly positive eigenvector.  The Perron coefficient of
every nonzero nonnegative initial vector is positive; dividing the spectral
expansion by $Z(t)$ proves the limit.
\end{proof}

The phrase \emph{projective semigroup owner} records the theorem content of
the original manuscript round.

\section{Complete finite-genome spectrum}

For $S\subseteq\{1,\ldots,L\}$ let
\[
 \phi_S(x)=2^{-L/2}(-1)^{\sum_{i\in S}x_i}.
\]
These Walsh characters form an orthonormal basis and
\begin{equation}\label{eq:mutation-spectrum}
 M_L\phi_S=d_{|S|}\phi_S,\qquad
 d_k=-\frac{2Uk}{L},\qquad
 \langle e_0,\phi_S\rangle=2^{-L/2}.
\end{equation}
Put $m_k=L!/[k!(L-k)!]$, $w_k=m_k/2^L$.

\begin{theorem}[Full multiplicity and secular factorization]\label{thm:spectrum}
The complete characteristic polynomial of $A_L$ is
\begin{align}
 \det(\lambda I-A_L)
 &=q_L(\lambda)\prod_{k=0}^L
       (\lambda-d_k)^{m_k-1},\label{eq:full-factor}\\
 q_L(\lambda)
 &=\prod_{k=0}^L(\lambda-d_k)
   -s\sum_{k=0}^Lw_k\prod_{j\ne k}(\lambda-d_j).
   \label{eq:secular-polynomial}
\end{align}
Thus $d_k$ is retained with multiplicity $m_k-1$ and the remaining $L+1$
eigenvalues are precisely the roots, away from the poles, of
\begin{equation}\label{eq:secular}
 1=s\sum_{k=0}^L\frac{w_k}{\lambda-d_k}.
\end{equation}
The multiplicities in~\eqref{eq:full-factor} sum to $2^L$.
\end{theorem}

\begin{proof}
Let $E_k=\operatorname{span}\{\phi_S:|S|=k\}$.  The subspace of $E_k$
whose Walsh coefficients sum to zero is orthogonal to $e_0$ and has dimension
$m_k-1$; $A_L$ acts there by $d_k$.  Its orthogonal complement over all $k$
has basis
\[
 u_k=m_k^{-1/2}\sum_{|S|=k}\phi_S,
 \qquad \langle e_0,u_k\rangle=\sqrt{w_k}.
\]
The restriction of $A_L$ is $D+svv^{\mathsf T}$, where
$D=\operatorname{diag}(d_0,\ldots,d_L)$ and $v_k=\sqrt{w_k}$.  The matrix
determinant lemma gives~\eqref{eq:secular-polynomial}; adjoining the retained
orthogonal spaces gives~\eqref{eq:full-factor}.
\end{proof}

\ifnum\CRevisionRound>0
\section{Strict interlacing and projective rate}

Order the mutation poles as $d_L<\cdots<d_1<d_0=0$ and set
$R(\lambda)=s\sum_kw_k/(\lambda-d_k)$.  On every pole-free interval,
\begin{equation}\label{eq:derivative}
 R'(\lambda)=-s\sum_{k=0}^L\frac{w_k}{(\lambda-d_k)^2}<0.
\end{equation}
At the left and right ends of each internal gap, $R$ has limits $+\infty$
and $-\infty$, respectively.  Above $d_0$, it decreases from $+\infty$ to
$0$; below $d_L$ it is negative.

\begin{theorem}[Interlacing and exact projective gap]\label{thm:gap}
The $L+1$ secular roots $\rho_0>\rho_1>\cdots>\rho_L$ are simple and obey
\begin{equation}\label{eq:interlace}
 \rho_0>d_0=0>\rho_1>d_1>\rho_2>\cdots>d_{L-1}>\rho_L>d_L.
\end{equation}
In particular, the two largest eigenvalues of the full $2^L$-dimensional
matrix are $\rho_0$ and $\rho_1$.  With
\begin{equation}\label{eq:gap}
 \Delta_L(U,s)=\rho_0-\rho_1,
\end{equation}
every simplex orbit satisfies
$p(t)-p_\infty=O(e^{-\Delta_Lt})$.  This exponent is attained whenever the
coefficient of the $\rho_1$ eigenvector in the projective expansion is
nonzero.
\end{theorem}

\begin{proof}
Equation~\eqref{eq:derivative} and the one-sided pole limits give exactly one
root in every internal gap, one above $d_0$, and none below $d_L$.  The
retained eigenvalue nearest the top is $d_1<\rho_1$, so $\rho_1$ is the true
second eigenvalue.  Divide the real-symmetric spectral expansion in
Theorem~\ref{thm:flow} by its Perron term.  All remaining exponents are at most
$\rho_1-\rho_0$; a nonzero $\rho_1$ coefficient gives the stated leading
term.
\end{proof}

The phrase \emph{strict secular interlacing} records the substantive Round-1
addition.
\fi

\ifnum\CRevisionRound>1
\section{Reducible faces and evidence boundary}

If $s=0$, selection disappears and $d_k$ has its full multiplicity $m_k$;
the mutation semigroup converges to the uniform law.  If $U=0$, writing
$a(t)=p_0(t)$ gives $\dot a=sa(1-a)$.  Initial data with $a(0)>0$ converge to
$e_0$, whereas the whole face $a(0)=0$ is stationary.  Thus no irreducible
Perron conclusion is smuggled onto that face.  For $L=1$, no $d_k$ is
retained and direct diagonalization gives
\begin{equation}\label{eq:l-one}
 \lambda_\pm=\frac{s-2U\pm\sqrt{s^2+4U^2}}2.
\end{equation}

The exact evidence enumerates rational parameter fixtures, reconstructs the
factorization independently, directly compares small hypercube matrices,
counts secular roots symbolically, checks Walsh eigenvectors and quotient
derivatives, and attacks signs, weights, multiplicities, scopes and parser
conventions.  These finite checks are regression receipts, not an extrapolation
proof.  At every fixed finite $L$ the strictly interlacing roots vary smoothly
in the positive chamber; this paper makes no singular infinite-genome error-
threshold claim.

\section{Ownership and Route-A boundary}

Crow--Kimura mutation--selection theory and the permutation-invariant linear
reduction are historical source owners.  Nearby workspace packages treat the
mutation-only Ehrenfest cube, Moran fixation, Wright--Fisher diffusion and
network SIS; none contains the single-peak rank-one full spectrum proved here.
The phrase \emph{finite-length boundary and route firewall} records the
substantive Round-2 addition.

The Hamming cube supplies neither rational-prime local data nor a logarithmic
prime clock.  This positive normalized flow has no isolated arithmetic
primitive ledger, and~\eqref{eq:full-factor} is a source characteristic
polynomial, not a target Euler factor or divisor.  Real symmetry gives only a
formal source-operator hint.  Under scope
\texttt{NO\_BAD\_EULER\_OR\_ROOT\_NUMBER}, the tuple is
\[
 (\mathrm{A0\_FAIL},\mathrm{A1\_FAIL},\mathrm{A2\_FAIL},
   \mathrm{A3\_FAIL},\mathrm{A4\_FORMAL\_HINT}).
\]
The verdict is \texttt{ROUTE\_A\_REJECTED}; Route B is locked.  No target
arithmetic datum, Euler factor, root number, automorphy, target functional
equation or zero match, or Hilbert--Polya operator is claimed.
\fi

\end{document}
